% Motor: XeLaTeX. Sesion 20 - material docente UTP.
% Misma estructura docente que la Sesion 19: teoria en dos partes y un lab.
\documentclass[aspectratio=169,10pt,c]{beamer}
% Reutiliza sin cambios la identidad visual UTP validada en la Sesion 09.
% Reutiliza sin cambios la identidad visual UTP validada en la Sesion 08.
% Reutiliza sin cambios la identidad visual UTP validada en la Sesion 05.
% La segunda ruta permite compartir los assets comunes sin duplicar binarios.
\input{../../Sesion05/slides/header.tex}
\graphicspath{{../../Sesion05/slides/assets/}}



% El mapa de regiones es la captura del curso que ya esta en la Sesion 17.
\graphicspath{{assets/}{../../Sesion17/slides/assets/}{../../Sesion05/slides/assets/}}
% Bloques de código con el mismo estilo que la Sesion 11.
\usepackage{listings}
\lstset{basicstyle=\ttfamily\small,backgroundcolor=\color{utpGrisClaro},
  frame=none,aboveskip=4pt,belowskip=4pt,showstringspaces=false,
  columns=fullflexible,keepspaces=true,breaklines=true,
  xleftmargin=5pt,xrightmargin=5pt,framesep=5pt}


\title{Integración de Servicios Cloud GCP}
\subtitle{Ruta Associate Cloud Engineer: rutas en VPC, Acceso privado a Google y Cloud NAT}
\date{Semana 10 - Sesión 20}
\author{Rel Guzman}

% Separador de tema: titulo grande, subtitulo y linea roja.
\newcommand{\separador}[2]{%
\begin{frame}[plain]
\vfill
\begin{center}
{\fontsize{29}{34}\selectfont\bfseries #1}\par
\vspace{2mm}
{\small #2}\par
\vspace{4mm}
{\color{utpRojo}\rule{8.5cm}{1.2pt}}
\end{center}
\vfill
\end{frame}}

\begin{document}

{\setbeamertemplate{footline}{}\maketitle}

% ------------------------------------------- Revision de la sesion anterior
\begin{frame}{Revisión de la sesión anterior (1/2): redes VPC}
\begin{block}{Lo que hicimos}
\small Vimos cómo se organiza Google Cloud y, en el lab \emph{Crea y configura redes de VPC}, creaste \textbf{tres redes} y probaste con \texttt{ping} qué máquinas se hablan.
\end{block}
\begin{columns}[T,onlytextwidth]
\column{.48\textwidth}
\begin{exampleblock}{Dos preguntas sobre cada recurso}
\small \textbf{¿De quién es?} Lo dice la jerarquía: organización, carpeta, proyecto.\par\smallskip
\textbf{¿Dónde está?} Lo dicen la región y la zona.
\end{exampleblock}
\column{.48\textwidth}
\begin{exampleblock}{El edificio}
\small La \textbf{red} es el edificio; la \textbf{subred}, un piso; el \textbf{gateway}, la puerta del piso.\par\smallskip
En Google Cloud eso se llama \textbf{red VPC}, y se arma sin cables.
\end{exampleblock}
\end{columns}
\par\vspace{1mm}
\begin{alertblock}{La idea que seguimos usando hoy}
\small Toda máquina virtual vive \textbf{dentro de una red}. Hoy vemos \textbf{por dónde viaja} lo que sale de ella.
\end{alertblock}
\end{frame}

\begin{frame}{Revisión de la sesión anterior (2/2): pregunta rápida}
\begin{block}{Lo que salió en el lab}
\small Desde una máquina hiciste \texttt{ping} a las otras tres, por dentro (IP interna) y por fuera (IP externa).
\end{block}
\renewcommand{\arraystretch}{1.3}
\centering\small
\begin{tabular}{@{}p{5.2cm}p{3.4cm}p{3.4cm}@{}}
\toprule
\textbf{La otra máquina está en\ldots} & \textbf{Por su IP interna} & \textbf{Por su IP externa}\\
\midrule
la \textbf{misma red}, en otra región & \textbf{Responde} & Responde\\
\textbf{otra red}, en la misma zona & \textbf{No llega} & Responde\\
\bottomrule
\end{tabular}
\par\vspace{3mm}
\begin{columns}[T,onlytextwidth]
\column{.48\textwidth}
\begin{exampleblock}{Lo que aprendimos}
\small Por dentro, lo que importa es estar en la \textbf{misma red}. Cada red VPC está aislada de las demás.
\end{exampleblock}
\column{.48\textwidth}
\begin{alertblock}{Pregunta rápida}
\small En el lab, todas las máquinas tenían IP externa. ¿Qué crees que pasa si a una \textbf{se la quitas}?
\end{alertblock}
\end{columns}
\end{frame}

\begin{frame}[c]
\begin{block}{Logro de la sesión}
Al finalizar, podrás explicar con tus palabras \textbf{qué es una red VPC} y \textbf{qué es una ruta}, y decir para qué sirven el \textbf{Acceso privado a Google} y \textbf{Cloud NAT}, haciendo que una máquina \textbf{sin dirección pública} llegue a Google y a Internet.
\end{block}
\end{frame}

\begin{frame}{Saberes previos}
\begin{itemize}\setlength\itemsep{10pt}
  \item Vas por la carretera a una ciudad que no conoces. ¿Cómo sabes por dónde ir en cada cruce?
  \item Desde tu celular puedes abrir cualquier página de Internet. ¿Puede alguien de Internet ``entrar'' a tu celular? ¿Por qué crees que no?
  \item En una oficina hay teléfonos que solo llaman a otros anexos. ¿Para qué serviría tener uno así?
\end{itemize}
\end{frame}

\begin{frame}[c]{Utilidad de esta sesión}
\begin{center}
\begin{minipage}{.86\textwidth}
{\large En el trabajo real, las máquinas más importantes \textbf{no se exponen a Internet}: una base de datos, un servidor interno. Pero igual necesitan \textbf{guardar archivos} y \textbf{actualizarse}. Hoy aprendes cómo se logra eso sin abrirles la puerta.}
\end{minipage}
\end{center}
\end{frame}

\begin{frame}{Plan de hoy: dos partes y un lab}
\renewcommand{\arraystretch}{1.35}
\centering\small
\begin{tabular}{@{}p{1.4cm}>{\raggedright\arraybackslash}p{6.4cm}>{\raggedright\arraybackslash}p{5.1cm}@{}}
\toprule
\textbf{Parte} & \textbf{Qué vemos} & \textbf{Del curso}\\
\midrule
1. Teoría & Repaso corto: qué es una red VPC y qué es una ruta & Módulo \emph{Redes virtuales}\\
2. Teoría & Máquinas sin dirección pública: cómo llegan a Google y a Internet & \emph{Acceso privado a Google y Cloud NAT}\\
\textbf{Lab} & Una VM privada que igual descarga y se actualiza & \emph{Implementa el Acceso privado a Google y Cloud NAT}\\
\bottomrule
\end{tabular}
\par\vspace{3mm}
\begin{alertblock}{Cómo lo trabajamos}
\small Vamos a lo \textbf{esencial}: una idea por diapositiva y una comparación de la vida diaria. Del lab vemos solo \textbf{lo más importante}; los pasos están en su guía.
\end{alertblock}
\end{frame}

% ------------------------------------------------ Parte 1: repaso de VPC
\separador{Parte 1: las redes VPC, en corto}{Qué es una red VPC y qué es una ruta}

\begin{frame}{¿Qué es una red VPC?}
\begin{block}{En una frase}
\small La \textbf{red privada} de tu proyecto dentro de Google Cloud: ahí se conectan tus máquinas virtuales. Se arma con clics o comandos, \textbf{sin cables}.
\end{block}
\begin{columns}[c,onlytextwidth]
\column{.54\textwidth}
\centering
\begin{tikzpicture}[font=\footnotesize]
  \draw[black!55,rounded corners=6pt] (-0.25,-0.25) rectangle (6.85,4.2);
  \node[anchor=north west,font=\small\bfseries] at (-0.15,4.15) {Proyecto};
  \draw[utpRojo,thick,rounded corners=6pt,fill=utpRojo!5] (0,0) rectangle (6.6,3.5);
  \node[anchor=north west,font=\small\bfseries] at (0.1,3.45) {Red (es mundial)};
  \draw[rounded corners=4pt,fill=white] (0.3,0.3) rectangle (3.1,2.7);
  \node[anchor=north west] at (0.35,2.65) {\textbf{Subred}};
  \node[anchor=north west] at (0.35,2.25) {región de EE.\,UU.};
  \node[draw,fill=utpRojo!20,rounded corners=2pt] at (1.05,0.9) {VM};
  \node[draw,fill=utpRojo!20,rounded corners=2pt] at (2.3,0.9) {VM};
  \draw[rounded corners=4pt,fill=white] (3.5,0.3) rectangle (6.3,2.7);
  \node[anchor=north west] at (3.55,2.65) {\textbf{Subred}};
  \node[anchor=north west] at (3.55,2.25) {región de Europa};
  \node[draw,fill=utpRojo!20,rounded corners=2pt] at (4.9,0.9) {VM};
\end{tikzpicture}
\column{.43\textwidth}
\begin{block}{De afuera hacia adentro}
\small
\textbf{Proyecto:} la caja donde vive todo.\par\smallskip
\textbf{Red:} abarca \textbf{todo el mundo}.\par\smallskip
\textbf{Subred:} un pedazo de la red en \textbf{una región}.\par\smallskip
\textbf{VM:} vive dentro de una subred.
\end{block}
\begin{exampleblock}{Una comparación}
\small Una \textbf{empresa} con una \textbf{sede} en cada ciudad y \textbf{empleados} en cada sede.
\end{exampleblock}
\end{columns}
\end{frame}

\begin{frame}{Las redes VPC en cinco ideas}
\renewcommand{\arraystretch}{1.4}
\centering\small
\begin{tabular}{@{}>{\raggedright\arraybackslash}p{3.7cm}>{\raggedright\arraybackslash}p{5.3cm}>{\raggedright\arraybackslash}p{3.9cm}@{}}
\toprule
\textbf{La idea} & \textbf{Qué significa} & \textbf{Se parece a}\\
\midrule
\textbf{La red es mundial} & Una sola red llega a todas las regiones & Una empresa con sedes en varios países\\
\textbf{La subred es de una región} & Cada una tiene su propio grupo de direcciones & La sede de una ciudad\\
\textbf{Misma red, se hablan} & Por dentro, aunque estén lejos. Entre redes distintas, no & Los anexos de una misma empresa\\
\textbf{Dos direcciones} & La IP interna la tienen todas. La externa es opcional & El anexo y el número público\\
\textbf{El firewall decide} & Lo que entra está bloqueado hasta que una regla lo permita & El portero con su lista\\
\bottomrule
\end{tabular}
\par\vspace{2mm}
\begin{alertblock}{La idea que abre la parte 2}
\small La IP externa es \textbf{opcional}. ¿Qué pasa con una máquina que \textbf{no la tiene}?
\end{alertblock}
\end{frame}

\begin{frame}{¿Qué es una ruta?}
\begin{center}
{\Large Una ruta es un \textbf{letrero}: ``para ir \textbf{allá}, ve por \textbf{aquí}''.}
\end{center}
\vspace{1mm}
\begin{center}
\begin{tikzpicture}[font=\small,
  caja/.style={draw=utpRojo,rounded corners=4pt,minimum width=3.2cm,minimum height=0.9cm,align=center}]
  \node[caja,fill=white] (v) at (0,0) {\textbf{Tu VM}};
  \node[caja,fill=utpRojo!10] (a) at (8,0.9) {\textbf{Otra VM de tu red}};
  \node[caja,fill=utpRojo!30] (b) at (8,-0.9) {\textbf{Internet}};
  \draw[-latex,thick] (v.east) -- node[above,sloped,font=\footnotesize] {por dentro} (a.west);
  \draw[-latex,thick] (v.east) -- node[below,sloped,font=\footnotesize] {hacia afuera} (b.west);
\end{tikzpicture}
\end{center}
\begin{columns}[T,onlytextwidth]
\column{.48\textwidth}
\begin{exampleblock}{Como en la carretera}
\small ``Arequipa $\rightarrow$''. El letrero no maneja por ti: solo te indica el camino.
\end{exampleblock}
\column{.48\textwidth}
\begin{alertblock}{Pero el letrero no basta}
\small Para salir a Internet, la VM necesita una \textbf{IP externa}. ¿Y si no tiene? Eso es la parte 2.
\end{alertblock}
\end{columns}
\par\vspace{2mm}
\centering\small Google pone estos dos letreros en toda red. \textbf{Tú no haces nada.}
\end{frame}

% --------------------------------------- Parte 2: Acceso privado y Cloud NAT
\separador{Parte 2: Acceso privado a Google y Cloud NAT}{Un ejemplo: vamos a desplegar una webapp}

\begin{frame}{Nuestro ejemplo: la webapp de una tienda}
\begin{center}
\begin{tikzpicture}[font=\small,
  caja/.style={draw=utpRojo,rounded corners=4pt,minimum width=3.2cm,minimum height=1.1cm,align=center}]
  \node[caja,fill=utpRojo!10] (s) at (0,0) {\textbf{Bucket}\\\footnotesize fotos de los productos};
  \node[caja,fill=white,very thick] (v) at (5.2,0) {\textbf{VM con la webapp}};
  \node[caja,fill=utpRojo!30] (i) at (10.4,0) {\textbf{Internet}\\\footnotesize actualizaciones};
  \draw[-latex,thick] (v) -- node[above,font=\footnotesize] {lee las fotos} (s);
  \draw[-latex,thick] (v) -- node[above,font=\footnotesize] {se actualiza} (i);
\end{tikzpicture}
\end{center}
\begin{columns}[T,onlytextwidth]
\column{.31\textwidth}
\begin{block}{Necesita 1}
\small Leer las \textbf{fotos} que guardamos en un bucket de \textbf{Cloud Storage}.
\end{block}
\column{.31\textwidth}
\begin{block}{Necesita 2}
\small Descargar \textbf{actualizaciones} y programas de Internet.
\end{block}
\column{.31\textwidth}
\begin{alertblock}{Y una regla}
\small \textbf{Nadie} de Internet debe poder entrar a esta VM.
\end{alertblock}
\end{columns}
\end{frame}

\begin{frame}{Por seguridad, la VM va sin IP externa}
\begin{center}
\begin{tikzpicture}[font=\small,
  caja/.style={draw=utpRojo,rounded corners=4pt,minimum width=3.2cm,minimum height=1.1cm,align=center},
  cortado/.style={-latex,thick,dashed,black!45}]
  \node[caja,fill=utpRojo!10] (s) at (0,0) {\textbf{Bucket}\\\footnotesize fotos};
  \node[caja,fill=white,very thick] (v) at (5.2,0) {\textbf{VM con la webapp}\\\footnotesize sin IP externa};
  \node[caja,fill=utpRojo!30] (i) at (10.4,0) {\textbf{Internet}};
  \draw[cortado] (v) -- node[font=\Large\bfseries,text=utpRojo] {$\times$} (s);
  \draw[cortado] ([yshift=2.2mm]v.east) -- node[above,font=\Large\bfseries,text=utpRojo] {$\times$} ([yshift=2.2mm]i.west);
  \draw[cortado] ([yshift=-2.2mm]i.west) -- node[below,font=\footnotesize,text=black] {nadie entra} ([yshift=-2.2mm]v.east);
\end{tikzpicture}
\end{center}
\begin{columns}[T,onlytextwidth]
\column{.48\textwidth}
\begin{exampleblock}{Lo bueno}
\small Desde Internet \textbf{nadie la ve}: nadie puede intentar entrar.
\end{exampleblock}
\column{.48\textwidth}
\begin{alertblock}{Lo malo}
\small Tampoco \textbf{sale}: no llega a las fotos ni a las actualizaciones. La webapp \textbf{no funciona}.
\end{alertblock}
\end{columns}
\par\vspace{1mm}
\begin{block}{Lo que buscamos}
\small Abrirle caminos de \textbf{salida} sin abrirle una puerta de \textbf{entrada}. Uno para las fotos y otro para Internet.
\end{block}
\end{frame}

\begin{frame}{Para las fotos: Acceso privado a Google}
\begin{center}
\begin{tikzpicture}[font=\small,
  caja/.style={draw=utpRojo,rounded corners=4pt,minimum width=3.2cm,minimum height=1.1cm,align=center}]
  \node[caja,fill=utpRojo!10] (s) at (0,0) {\textbf{Bucket}\\\footnotesize fotos};
  \node[caja,fill=white,very thick] (v) at (7,0) {\textbf{VM con la webapp}\\\footnotesize sin IP externa};
  \draw[-latex,thick] (v) -- node[above,font=\footnotesize] {por dentro de Google} (s);
\end{tikzpicture}
\end{center}
\begin{columns}[T,onlytextwidth]
\column{.48\textwidth}
\begin{block}{Qué es}
\small Un \textbf{interruptor} de la \textbf{subred}. Lo enciendes y sus VMs llegan a los servicios de Google, como Cloud Storage.
\end{block}
\begin{alertblock}{Solo llega a Google}
\small Las actualizaciones de Internet \textbf{siguen sin llegar}.
\end{alertblock}
\column{.48\textwidth}
\begin{exampleblock}{Como en una empresa}
\small Un \textbf{pasillo interno} hasta el almacén: llegas sin salir a la calle.
\end{exampleblock}
\begin{block}{Dónde se enciende}
\small En la \textbf{subred}, no en la VM. Vale para todas las VMs de esa subred.
\end{block}
\end{columns}
\end{frame}

\begin{frame}{Para las actualizaciones: Cloud NAT}
\begin{center}
\begin{tikzpicture}[font=\small,
  caja/.style={draw=utpRojo,rounded corners=4pt,minimum width=2.9cm,minimum height=1.1cm,align=center}]
  \node[caja,fill=white,very thick] (v) at (0,0) {\textbf{VM con la webapp}\\\footnotesize sin IP externa};
  \node[caja,fill=utpRojo!30] (n) at (4.6,0) {\textbf{Cloud NAT}\\\footnotesize le presta una IP};
  \node[caja,fill=utpRojo!10] (i) at (10.6,0) {\textbf{Internet}};
  \draw[-latex,thick] ([yshift=2.2mm]v.east) -- node[above,font=\footnotesize] {sale} ([yshift=2.2mm]n.west);
  \draw[-latex,thick] ([yshift=2.2mm]n.east) -- node[above,font=\footnotesize] {sale con la IP prestada} ([yshift=2.2mm]i.west);
  \draw[-latex,thick,dashed] ([yshift=-2.2mm]i.west) -- node[below,font=\footnotesize] {nadie entra por su cuenta} ([yshift=-2.2mm]n.east);
\end{tikzpicture}
\end{center}
\begin{columns}[T,onlytextwidth]
\column{.48\textwidth}
\begin{block}{Qué es}
\small Deja a la VM \textbf{salir} a Internet. La respuesta vuelve, pero nadie de afuera puede \textbf{empezar} una conexión hacia ella.
\end{block}
\column{.48\textwidth}
\begin{exampleblock}{Ya lo usas en casa}
\small Tu \textbf{router wifi}: tu celular navega y nadie de Internet entra a él.
\end{exampleblock}
\end{columns}
\par\vspace{1mm}
\begin{block}{¿Y NAT?}
\small \emph{Traducción de direcciones}: al salir, cambia la IP interna de la VM por una pública.
\end{block}
\end{frame}

\begin{frame}{La webapp completa}
\begin{center}
\begin{tikzpicture}[font=\small,
  caja/.style={draw=utpRojo,rounded corners=4pt,minimum width=3.0cm,minimum height=1.0cm,align=center}]
  \node[caja,fill=utpRojo!10] (s) at (0,0) {\textbf{Bucket}\\\footnotesize fotos};
  \node[caja,fill=white,very thick] (v) at (5.2,0) {\textbf{VM con la webapp}\\\footnotesize sin IP externa};
  \node[caja,fill=utpRojo!30] (i) at (10.4,0) {\textbf{Internet}};
  \draw[-latex,thick] (v) -- node[above,font=\footnotesize] {Acceso privado} (s);
  \draw[-latex,thick] ([yshift=2.2mm]v.east) -- node[above,font=\footnotesize] {Cloud NAT} ([yshift=2.2mm]i.west);
  \draw[-latex,thick,dashed,black!45] ([yshift=-2.2mm]i.west) -- node[below,font=\footnotesize,text=black] {nadie entra} ([yshift=-2.2mm]v.east);
\end{tikzpicture}
\end{center}
\renewcommand{\arraystretch}{1.35}
\centering\small
\begin{tabular}{@{}>{\raggedright\arraybackslash}p{4.6cm}>{\raggedright\arraybackslash}p{4.3cm}p{3.6cm}@{}}
\toprule
\textbf{La webapp necesita\ldots} & \textbf{Enciendes} & \textbf{Dónde}\\
\midrule
Leer las fotos del bucket & \textbf{Acceso privado a Google} & en la subred\\
Descargar actualizaciones & \textbf{Cloud NAT} & en la red y la región\\
\bottomrule
\end{tabular}
\par\vspace{2mm}
\begin{alertblock}{El resultado}
\small La webapp \textbf{trabaja con normalidad} y sigue \textbf{invisible} desde Internet. En el lab haces justo esto, en este orden.
\end{alertblock}
\end{frame}

% ------------------------------------------------------------------ Lab
\separador{El lab: lo más importante}{\emph{Implementa el Acceso privado a Google y Cloud NAT}}

\begin{frame}{El lab en una mirada}
\begin{block}{De qué trata}
\small Creas una máquina \textbf{sin IP externa} y le vas abriendo caminos de salida, uno por uno.
\end{block}
\begin{center}
\begin{tikzpicture}[font=\small,
  caja/.style={draw=utpRojo,rounded corners=4pt,minimum width=3.5cm,minimum height=1.15cm,align=center}]
  \node[caja,fill=white] (a) at (0,0) {\textbf{1. VM sin IP externa}\\\footnotesize no llega a ningún lado};
  \node[caja,fill=utpRojo!15] (b) at (4.7,0) {\textbf{2. + Acceso privado}\\\footnotesize ya llega a Google};
  \node[caja,fill=utpRojo!30] (c) at (9.4,0) {\textbf{3. + Cloud NAT}\\\footnotesize ya llega a Internet};
  \draw[-latex,thick] (a) -- (b);
  \draw[-latex,thick] (b) -- (c);
\end{tikzpicture}
\end{center}
\begin{columns}[T,onlytextwidth]
\column{.48\textwidth}
\begin{block}{Las dos pruebas del lab}
\small \textbf{¿Llego a Google?} Traer una imagen de un bucket, como las fotos de la tienda.\par\smallskip
\textbf{¿Llego a Internet?} Pedir actualizaciones con \texttt{sudo apt-get update}, como la webapp.
\end{block}
\column{.48\textwidth}
\begin{exampleblock}{El ritmo del lab}
\small \textbf{No funciona}, enciendes algo, \textbf{ahora funciona}. Así, dos veces.
\end{exampleblock}
\end{columns}
\par\vspace{1mm}
\begin{alertblock}{No te asustes}
\small Que la prueba falle al inicio \textbf{es parte del lab}: no es un error tuyo. Si se queda esperando, corta con \texttt{Ctrl + C}.
\end{alertblock}
\end{frame}

\begin{frame}{Lo más importante del lab: antes y después}
\begin{block}{La misma máquina, tres momentos}
\small \texttt{vm-internal} nunca recibe una IP externa. Lo único que cambia es lo que vas \textbf{encendiendo} a su alrededor.
\end{block}
\renewcommand{\arraystretch}{1.35}
\centering\small
\begin{tabular}{@{}p{4.9cm}p{2.3cm}p{2.6cm}p{2.6cm}@{}}
\toprule
\textbf{Prueba dentro de la VM} & \textbf{Al inicio} & \textbf{+ Acceso privado} & \textbf{+ Cloud NAT}\\
\midrule
Traer la imagen del bucket & No llega & \textbf{Funciona} & Funciona\\
\texttt{sudo apt-get update} & No llega & Solo lo de Google & \textbf{Funciona}\\
Que alguien de Internet entre & No & No & \textbf{No}\\
\bottomrule
\end{tabular}
\par\vspace{2mm}
\begin{columns}[T,onlytextwidth]
\column{.48\textwidth}
\begin{exampleblock}{Qué significa}
\small Cada servicio abre \textbf{un camino de salida} distinto: uno hacia Google, el otro hacia Internet.
\end{exampleblock}
\column{.48\textwidth}
\begin{alertblock}{La conclusión del lab}
\small La máquina \textbf{trabaja con normalidad} y sigue \textbf{invisible} desde Internet.
\end{alertblock}
\end{columns}
\end{frame}

\begin{frame}{El lab: tres ideas para no perderte}
\renewcommand{\arraystretch}{1.45}
\centering\small
\begin{tabular}{@{}>{\raggedright\arraybackslash}p{4.3cm}>{\raggedright\arraybackslash}p{4.6cm}>{\raggedright\arraybackslash}p{4.0cm}@{}}
\toprule
\textbf{La idea} & \textbf{Dónde la ves en el lab} & \textbf{Si la olvidas}\\
\midrule
\textbf{La VM nace sin IP externa} & Al crearla, en IP externa eliges \textbf{Ninguna} & El lab pierde sentido: todo funciona desde el inicio\\
\textbf{El Acceso privado es de la subred} & Editas la \textbf{subred}, no la máquina & Lo buscas en la VM y no aparece\\
\textbf{Cloud NAT no es instantáneo} & Tarda hasta \textbf{3 minutos} en hacer efecto & Crees que falló: espera y repite\\
\bottomrule
\end{tabular}
\par\vspace{3mm}
\begin{columns}[T,onlytextwidth]
\column{.48\textwidth}
\begin{exampleblock}{¿Y cómo entras a la VM?}
\small No tiene IP externa, así que Google te abre un \textbf{túnel seguro} (IAP). El comando ya viene escrito en el lab.
\end{exampleblock}
\column{.48\textwidth}
\begin{alertblock}{Si algo falla}
\small Fíjate \textbf{dónde estás}: Cloud Shell y la VM son dos lugares distintos. Lo que guardaste en uno no existe en el otro.
\end{alertblock}
\end{columns}
\end{frame}

% ------------------------------------------------------------ Practica
\separador{Práctica}{El lab real}

\begin{frame}{Práctica: una VM sin dirección pública \textnormal{\small(en clase + tarea)}}
\begin{block}{En clase}
\small
1. Haz el lab real: \emph{Implementa el Acceso privado a Google y Cloud NAT}.\par\smallskip
2. En tu cuaderno, anota cada prueba del lab: \textbf{¿falló o funcionó?} y \textbf{¿qué acababas de encender?}\par\smallskip
3. Al terminar, compara tus notas con la tabla de \emph{antes y después}.
\end{block}
\begin{columns}[T,onlytextwidth]
\column{.48\textwidth}
\begin{exampleblock}{De tarea}
\small Terminar el lab si no alcanzó el tiempo en clase.
\end{exampleblock}
\column{.48\textwidth}
\begin{alertblock}{Entregable}
\small Captura del lab con sus tareas en verde y tus notas de las pruebas.
\end{alertblock}
\end{columns}
\end{frame}

\begin{frame}{Cierre: preguntas para comprobar}
\begin{block}{Responde con tus palabras}
\small
1. ¿Qué es una red VPC? ¿En qué se diferencia de una subred?\par\medskip
2. ¿Qué dice una ruta? ¿Qué dos tipos de ruta trae toda red VPC?\par\medskip
3. La ruta hacia Internet ya existe. ¿Basta para que una máquina sin IP externa salga a Internet?\par\medskip
4. ¿Por qué querrías una máquina sin IP externa?\par\medskip
5. La VM de tu webapp no tiene IP externa y necesita leer fotos de un bucket y descargar actualizaciones. ¿Qué enciendes para cada cosa?\par\medskip
6. Con Cloud NAT encendido, ¿puede alguien de Internet entrar a la máquina? ¿Por qué?
\end{block}
\begin{alertblock}{Lo que viene}
\small Las \textbf{máquinas virtuales} por dentro: tipos de máquina, discos e imágenes.
\end{alertblock}
\end{frame}

\begin{frame}[plain]
\begin{tikzpicture}[remember picture,overlay]
  \node[anchor=center,inner sep=0pt] at (current page.center)
    {\includegraphics[width=\paperwidth,height=\paperheight]{gracias-final.pdf}};
\end{tikzpicture}
\end{frame}

\begin{frame}{Referencias}
\scriptsize
\begin{itemize}\setlength\itemsep{3pt}
  \item Google Skills: ruta \emph{Certificación de Associate Cloud Engineer}, \texttt{skills.google/paths/11}.
  \item Google Skills: \emph{Infraestructura esencial de Google Cloud: conceptos básicos}, \texttt{skills.google/course\_templates/50}, módulo \emph{Redes virtuales}.
  \item Google Skills: lab \emph{Implementa el Acceso privado a Google y Cloud NAT}, \texttt{skills.google/catalog\_lab/2526}.
  \item Google Cloud: \emph{VPC networks}, \texttt{cloud.google.com/vpc/docs/vpc}.
  \item Google Cloud: \emph{Routes}, \texttt{cloud.google.com/vpc/docs/routes}.
  \item Google Cloud: \emph{Private Google Access}, \texttt{cloud.google.com/vpc/docs/private-google-access}.
  \item Google Cloud: \emph{Cloud NAT overview}, \texttt{cloud.google.com/nat/docs/overview}.
  \item Material de repaso: \texttt{Sesion19/slides/slides.tex}.
  \item Las respuestas de todas las preguntas de esta sesión están en \texttt{respuestas.md}.
\end{itemize}
\end{frame}

\end{document}
